\section{Consensus Engine Binding}
\label{sec:binding}

\subsection{The Lux consensus engine}

The Lux consensus engine \cite{consensusrepo} is a unified driver for
two modes: linear-chain (Nova) and DAG (Nebula). Both modes follow
the same per-round protocol family --- Photon (committee selection),
Wave (threshold voting + Fast Probabilistic Consensus), Focus
(consecutive-success counter), Prism (DAG geometry / cuts), Flare
(cert / skip via $2f+1$ quorum), Horizon (DAG safe-prefix), Ray
(linear finality driver), Field (DAG finality driver). The full
package layout is documented at \texttt{luxfi/consensus/protocol/}.

Quasar lives at \texttt{protocol/quasar/} and is the cert-producing
layer that the consensus driver invokes once a $2f+1$ quorum is
observed. The driver does not know what kind of cert Quasar produces;
it only knows that Quasar emits a \texttt{QuasarCert} byte string
that it forwards to the block's \texttt{Cert} field. Verification on
incoming blocks is symmetrical: the driver passes the cert bytes to
\texttt{quasar.Verify}; the engine treats the result as a black-box
boolean.

This is the orthogonality discipline: the consensus engine handles
ordering, voting, and reachability; Quasar handles cryptographic
finality; the QuasarCert is the only shared interface between them.

\subsection{Per-round witness pattern}

A consensus round at height $h$ proceeds as follows. The committee
$C_h$ is selected via Photon's $K$-of-$N$ scheme; the proposer
broadcasts the block $B_h$ over the ZAP wire protocol
\cite{lp020}; voting proceeds via Wave until a $2f+1$ quorum is
observed; the engine invokes
\texttt{quasar.Sign(committee, block, profile)} to produce
$\mathcal{C}_h$.

Inside \texttt{quasar.Sign}, the legs are signed in parallel:

\begin{itemize}[leftmargin=2em,nosep]
\item The BLS leg is the aggregate of per-validator BLS signatures
      on the round digest $\mu$. Each validator signs locally and
      broadcasts; the engine collects $\geq \tau$ signatures and
      aggregates by point addition (Section~\ref{sec:prelim}).
\item The Pulsar leg is produced by the Pulsar two-round threshold
      protocol \cite{pulsarrepo}: Round-1 commitments
      $\mathbf{w}_i = \mathbf{A}\mathbf{y}_i$ are broadcast; the
      aggregator collects them, computes the challenge $c =
      H_{\mathsf{Puls}}(\mu \,\|\, \sum_i \mathbf{w}_i)$, and
      forwards to Round-2; validators emit $\mathbf{z}_i$; the
      aggregator combines via Lagrange interpolation and emits the
      byte-equal FIPS 204 signature.
\item The Corona leg (on Aurora / Polaris) follows the same
      two-round structure on the R-LWE ring
      \cite{coronarepo,boschini2024corona}.
\item The Magnetar leg (on Polaris) follows the v0.1 reveal-and-aggregate
      protocol \cite{magnetarrepo}: validators emit Shamir shares of
      the master seed; the aggregator reconstructs $S$, calls
      \texttt{slhdsa.SignDeterministic}, and zeroizes $S$.
\item The Z-Chain leg is produced asynchronously by Z-Chain validators:
      per-validator ML-DSA-65 signatures are streamed to Z-Chain,
      which produces a Groth16 proof of valid signature quorum over a
      window of $W \approx 1024$ rounds. The proof is emitted once
      per window, not per block; for blocks within a window without
      a fresh proof, the cert carries the most recent valid proof
      bound to the same epoch.
\end{itemize}

The legs are signed in parallel because they are independent
cryptographic operations on the same round digest $\mu$. The cert is
assembled by concatenating the leg bytes per the wire format of
Section~\ref{sec:wire}, and is broadcast as part of the block's
\texttt{Cert} field.

\subsection{Verification flow}

On the receiving side, the consensus engine invokes
\texttt{quasar.Verify(cert, block, validators, profile)} which
performs the predicate of Definition~\ref{def:verify}. The verify
function is stateless --- it depends only on its arguments --- and
returns a boolean. The engine accepts the block iff verify returns
true; rejection causes the block to be dropped from the local DAG
or chain.

Because each leg has its own verifier, the legs may be verified in
parallel. The Lux \texttt{quasar.Verify} implementation dispatches
each leg to its primitive verifier as a goroutine, with a fast-fail:
if any leg returns false, the others are cancelled. This achieves
the lowest verify latency (bounded by the slowest leg, not the sum).

\subsection{Quasar in the linear-chain and DAG drivers}

In the linear-chain driver (\texttt{Nova}), Quasar is invoked once
per finalized block. The cert is stored on the block and propagated
to peers via ZAP.

In the DAG driver (\texttt{Nebula}), Quasar is invoked once per
finalized vertex along the DAG's safe-prefix commitment line
\cite{lp020}. The cert is stored on the vertex and the DAG store
(\texttt{core/dag/}) indexes certs by vertex ID for light-client
querying.

Both drivers use the \emph{same} \texttt{quasar.QuasarCert} byte
structure --- the wire format is independent of the consensus mode.
This is critical because Lux supports mixed-mode subnet topologies:
a chain may run linear-chain consensus while its parent runs DAG
consensus, and certs flow between them via Warp messaging.

\subsection{Warp messaging carries cert finality}

Cross-chain messaging on Lux uses the Warp protocol
\cite{quasarcertsoundness,lp075}: a message from chain $c_1$ to
chain $c_2$ is signed by $c_1$'s validators and is verified by
$c_2$'s validators using $c_1$'s known public keys.

A Warp message carries the source chain's QuasarCert as the
finality witness. The destination chain runs
\texttt{quasar.Verify(cert, block, validators, profile)} against the
source chain's profile and validator set, then trusts the message
as finalized.

This means the QuasarCert is the universal finality envelope: every
cross-chain interaction in the Lux ecosystem carries a verifiable
PQ-finality witness. The wire-format stability commitment in
Section~\ref{sec:wire} is what makes this work: a single
\texttt{QuasarCert} can be verified against any Lux chain's
validator set, profile, and primitives.

\subsection{Z-Chain rollup async dependency}

The Z-Chain Groth16 leg has a unique structure: it is not produced
per-block on the source chain; it is produced asynchronously by
Z-Chain validators over a window. This means the $\mathrm{ZK}$ leg
of $\mathcal{C}_h$ may carry a Groth16 proof produced at height
$h - k$ for some $k \in [0, W)$, where $W$ is the rollup window
size.

The freshness invariant is that the proof must be bound to the same
\emph{epoch} as the block; an epoch boundary forces a fresh proof.
This is enforced by the round-digest construction
(Section~\ref{sec:prelim}): $\mu$ includes the epoch $e$, so a
proof bound to epoch $e$ cannot be replayed for epoch $e' \neq e$.

Within an epoch, the proof's freshness window provides bandwidth
amortization: a single Groth16 proof can cover up to $W \approx 1024$
blocks before a new proof is required. This amortizes the
${\sim}400$ ms CPU prover time (or ${\sim}5$--$15$ ms GPU prover time;
see Appendix B of \cite{quasarcertsoundness}) across the window.

The trade-off: a block at height $h$ may carry a proof produced for
the validator-set state at height $h - k$. If a validator changes
keys mid-epoch, the proof remains valid against the older root, but
new validators may not appear in the proof. This is acceptable
because the validator set is stable within an epoch (validator-set
rotations happen at epoch boundaries by construction).

\subsection{Dynamic resharing}

The Pulsar / Corona / Magnetar group public keys must remain stable
across validator-set rotations, otherwise a key rotation would
invalidate the cert format. This is achieved via verifiable
secret resharing \cite{luxlssmpc}: when the validator set rotates
from $\mathcal{V}_e$ to $\mathcal{V}_{e+1}$, each retiring validator
shares its threshold secret with the incoming set under the same
group public key.

The Lux Secret Sharing (LSS) framework \cite{luxlssmpc} is the
generic adapter for this lifecycle. Pulsar instantiates LSS with
the Pulsar key-era / resharing / activation-cert lifecycle
\cite{pulsarrepo}; Corona will do the same once its Tier A
documentation completes (current status Tier A docs landed
\cite{coronarepo}); Magnetar's reshare semantics under the v0.1
reveal-and-aggregate construction are degenerate (because the
master seed is the durable secret, not the per-party shares;
resharing is a fresh DKG).

The key observation is that resharing is \emph{orthogonal} to the
QuasarCert format: the cert carries no resharing transcript; the
resharing is governed by the LSS lifecycle and committed to the
\texttt{validator\_root}$_e$ Merkle root at epoch boundaries. The
cert sees only the post-resharing public keys, never the resharing
process itself.
